\section{``Principles of Econometrics'' Handbook}
Each \gretlpy distribution includes a set of scripts
that demonstrate how to solve examples used in the renowned
econometrics handbook \cite{HillGriffithsLim_principles_2012}.
All files are placed in \texttt{Principles\_of\_Econometrics\_4th\_edition},
a subdirectory of \texttt{examples}. Its contents are as follows:
\begin{description}
	\item[\bftt{data}] a subdirectory containing all
		datasets used in the examples (with special thanks
		to Lee Adkins).
	\item[\bftt{functions}] a subdirectory containing the
		\texttt{hausman\_taylor.gfn} package for performing
		Hausman-Taylor estimation.\newline
		\textbf{Note:} this file may be removed in future \gretlpy
		releases, since the package is officially available
		through SourceForge.
	\item[\bftt{scripts}] scripts demonstrating examples from
		the handbook chapters: \bftt{chapter\_5.py}, \bftt{chapter\_6.py},
		\bftt{chapter\_7.py}, \bftt{chapter\_8.py}, \bftt{chapter\_9.py},
		\bftt{chapter\_10.py}, \bftt{chapter\_11.py}, \bftt{chapter\_12.py},
		\bftt{chapter\_13.py}, \bftt{chapter\_14.py}, \bftt{chapter\_15.py},
		and \bftt{chapter\_16.py}.
\end{description}
To print out the contents of \bftt{chapter\_13.py}, the user
should call the following snippet:
\begin{lstlisting}[style=pythonstylefoot]
import gretl

gretl.print_module_source("chapter_13")
\end{lstlisting}\index{print\_module\_source()}

Figure~\ref{fig:poe4_chapter_13} shows the output
of the above snippet.

\begin{figure}
	\centering
	\includegraphics[width=.9\columnwidth,keepaspectratio]{poe4_example_chapter_13}
	\caption{Output of \texttt{chapter\_13.py} as displayed by \pythcode{gretl.print\_module\_source()}.}
	\label{fig:poe4_chapter_13}
\end{figure}
